% =====================================================================
%  Lecture 03 -- 2G/3G Core Networks                      (English edition)
% =====================================================================
\documentclass[11pt,a4paper]{article}
\providecommand\ACNroot{../..}
\input{\ACNroot/templates/lecture-en}

\ACNset{number}{03}
\ACNset{title}{2G/3G Core Networks}
\ACNset{shorttitle}{2G/3G Core Networks}
\ACNset{titleother}{شبکه‌های هستهٔ نسل دوم و سوم}
\ACNset{part}{Part I -- 2G/3G foundations (Lectures 1--4)}
\ACNset{time}{6--8 hours (3 h reading, 3--5 h exercises)}
\ACNset{level}{Intermediate}
\ACNset{prereq}{Lectures 1--2; spread spectrum basics; decibels}
\ACNset{gen}{3}

\begin{document}
\ACNcover

% ---------------------------------------------------------------------
\section{Learning objectives}
After studying this lecture you should be able to:
\begin{objectives}
  \item \textbf{Explain} the requirements that drove 3G (IMT-2000) and why a new radio
        interface, WCDMA, was needed.
  \item \textbf{Draw} the UMTS architecture -- UE, NodeB, RNC, CS and PS core domains -- and
        name the interfaces Uu, Iub, Iur, Iu-CS and Iu-PS.
  \item \textbf{Compute} processing gain, the uplink load factor and the noise rise of a WCDMA
        cell, and \textbf{estimate} its voice pole capacity.
  \item \textbf{Describe} the HSPA evolution (Releases 5--8) and \textbf{derive} the HSDPA peak rate
        from codes, spreading factor and modulation.
  \item \textbf{Trace} an RRC connection and radio access bearer (RAB) set-up with RANAP.
  \item \textbf{Explain} the role of IMS and \textbf{evaluate} why LTE removed the CS/PS split.
\end{objectives}

% ---------------------------------------------------------------------
\section{Prerequisite knowledge}
This lecture builds directly on the GSM/GPRS architecture of Lecture~1 and the
circuit-switched core of Lecture~2. You should also recall that a spread-spectrum signal
occupies a bandwidth much larger than its data rate, and that a receiver correlating
with the right code recovers the signal while other codes look like noise.

\paragraph{Key terms} The vocabulary of this lecture with its Persian equivalents:
\begin{keyterms}
UMTS & \fa{سامانهٔ جهانی مخابرات موبایل} & The 3GPP 3G system: WCDMA radio + evolved GSM core. \\
WCDMA & \fa{دسترسی چندگانهٔ تقسیم کد پهن‌باند} & 3G radio interface: \qty{5}{\mega\hertz} carrier, 3.84\,Mchip/s. \\
NodeB & \fa{نود بی} & 3G base station: radio transmission, spreading, inner-loop power control. \\
RNC & \fa{کنترل‌گر شبکهٔ رادیویی} & Controls NodeBs: radio resources, handover, ciphering. \\
Spreading factor (SF) & \fa{ضریب گسترش} & Chips per symbol; trades rate against protection. \\
Processing gain & \fa{بهرهٔ پردازش} & $W/R$: how much despreading lifts the wanted signal. \\
Soft handover & \fa{تحویل نرم} & UE connected to several cells at once (make-before-break). \\
Noise rise & \fa{افزایش نویز} & Increase of uplink interference over thermal noise. \\
RAB & \fa{حامل دسترسی رادیویی} & Bearer between UE and core network, set up by RANAP. \\
HSPA & \fa{دسترسی بستهٔ پرسرعت} & Shared-channel evolution of WCDMA (HSDPA + HSUPA). \\
IMS & \fa{زیرسامانهٔ چندرسانه‌ای IP} & SIP-based control of IP multimedia services. \\
\end{keyterms}

% ---------------------------------------------------------------------
\section{Introduction and motivation}
By the late 1990s GSM had proved that mobile telephony could be global, but it could not
become a data network. EDGE squeezed about 59\,kbit/s out of a timeslot, and one 200\,kHz
carrier can never give a single user megabits per second. The ITU's IMT-2000 programme
asked for up to 2\,Mbit/s indoors, 384\,kbit/s for pedestrians and 144\,kbit/s for vehicles,
flexible multiplexing of services with different rates on one connection, and global
roaming.

3GPP's answer, UMTS, combined two decisions. On the radio side it adopted a completely new
interface, Wideband CDMA, with a \qty{5}{\mega\hertz} carrier. On the network side it did the
opposite: it \emph{reused} the GSM/GPRS core -- MSC, HLR, SGSN, GGSN -- and connected it to the
new radio network through a new interface, Iu. This lecture studies that hybrid: a new radio
access network attached to an evolved 2G core with two parallel domains, and the strains
that eventually led LTE to abandon the circuit-switched domain altogether.

% ---------------------------------------------------------------------
\section{Fundamental concepts}

\subsection{Code division multiple access}
In WCDMA all users of a cell transmit at the same time on the same \qty{5}{\mega\hertz} carrier.
Each user's symbols are multiplied by a code running at the chip rate
$W = 3.84$\,Mchip/s. The receiver multiplies again by the same code (despreading): the wanted
signal is concentrated back into its narrow data bandwidth, while signals with other codes
remain spread and act like noise.

\begin{definitionbox}{Spreading factor and processing gain}
The spreading factor is the number of chips per symbol, $\mathrm{SF} = W/R_{\mathrm{s}}$. For a
service with information rate $R$, the processing gain is $G_{\mathrm{p}} = W/R$. It is the
factor by which despreading increases the signal-to-interference ratio.
\end{definitionbox}

Two families of codes are used. \emph{Channelisation codes} (orthogonal variable spreading
factor, OVSF) separate the channels sent by one transmitter; \emph{scrambling codes} separate
transmitters (cells in the downlink, UEs in the uplink). OVSF codes form a tree
(\cref{fig:ovsf}); a code can be used only if no ancestor or descendant is in use, so assigning
a short (low-SF, high-rate) code blocks a whole branch.

\subsection{Interference-limited capacity}
Because everyone shares one carrier, WCDMA capacity is limited by interference, not by a fixed
number of channels: every additional user raises the interference seen by all the others. Two
consequences follow. First, \emph{power control} is essential: an uplink user close to the
NodeB must not drown distant users (the near--far problem), so the NodeB adjusts each UE's power
1500 times per second. Second, capacity is \emph{soft}: the operator can trade a little quality
for more users, and the cell coverage shrinks as the load grows (cell breathing).

\subsection{Soft handover}
Since neighbouring cells use the same frequency, a UE near a cell border can be connected to two
or three cells simultaneously. In the uplink the RNC selects the best frame; in the downlink the
UE combines the signals of all cells in its active set. Soft handover avoids the gap of a hard
handover and provides macro-diversity, at the cost of extra NodeB and Iub resources. It is
analysed in Lecture~4.

% ---------------------------------------------------------------------
\section{Detailed technical explanation}

\subsection{UTRAN: NodeB and RNC}
The UMTS Terrestrial Radio Access Network (UTRAN) \cite{ts25401} has two node types.
The \textbf{NodeB} performs the physical layer: channel coding, spreading, modulation, RF and
the fast inner-loop power control. The \textbf{RNC} (Radio Network Controller) is much more
powerful than a GSM BSC. It terminates the radio protocols RLC, MAC (for dedicated channels)
and RRC, performs admission control, handover decisions, outer-loop power control, macro-diversity
combining, and -- unlike GSM, where ciphering ended at the BTS -- ciphering and integrity
protection. An RNC plays different roles for a given UE:
\begin{itemize}
  \item \textbf{SRNC} (serving RNC): terminates the Iu connection and the RRC signalling of that UE;
  \item \textbf{DRNC} (drift RNC): lends radio resources of its cells when the UE is in soft
        handover with cells of another RNC; data passes over \iface{Iur} to the SRNC;
  \item \textbf{CRNC} (controlling RNC): the RNC that controls a given NodeB's resources.
\end{itemize}

\subsection{Interfaces}
\begin{table}[htbp]
\centering\small
\caption{UMTS interfaces.}\label{tab:if}
\rowcolors{2}{acnmist}{white}
\begin{tabularx}{\linewidth}{@{}l l l Y@{}}
\toprule
\rowcolor{white}Interface & Between & Control protocol & Notes \\
\midrule
Uu    & UE -- NodeB/RNC & RRC              & WCDMA radio interface \\
Iub   & NodeB -- RNC    & NBAP             & Transport originally ATM, later IP \\
Iur   & RNC -- RNC      & RNSAP            & Soft handover across RNCs; no GSM equivalent \\
Iu-CS & RNC -- MSC      & RANAP            & Replaces the A interface; AMR speech \\
Iu-PS & RNC -- SGSN     & RANAP            & User plane in GTP-U; replaces Gb \\
\bottomrule
\end{tabularx}
\end{table}

A key design change is that the Iu-PS user plane uses GTP-U all the way from the RNC, so the
SGSN no longer handles LLC/SNDCP towards the terminal; the radio-specific protocols end in the
RNC. RANAP \cite{ts25413} carries the signalling for RAB management, relocation of the SRNC,
paging and security mode control.

\subsection{The core network: CS and PS domains}
The UMTS core keeps the two domains of GSM/GPRS:
\begin{itemize}
  \item The \textbf{circuit-switched (CS) domain} (MSC/VLR, GMSC) serves voice and video calls
        with Erlang-dimensioned circuits. From Release~4 the MSC is split into an \emph{MSC Server}
        (call control and mobility, the control plane) and a \emph{Media Gateway} (MGW, the user
        plane that switches and transcodes the speech). This allows the CS core transport to run
        over IP or ATM -- an early example of control/user plane separation.
  \item The \textbf{packet-switched (PS) domain} (SGSN, GGSN) serves IP traffic with PDP
        contexts, now with QoS classes (conversational, streaming, interactive, background).
\end{itemize}
The HLR/AuC is shared; the AuC now produces \emph{quintets}
$(\text{RAND},\text{XRES},\text{CK},\text{IK},\text{AUTN})$ for mutual authentication
(\emph{UMTS AKA}, \cite{ts33102}). AUTN lets the USIM verify that the challenge came from its
home network and is fresh; IK protects signalling integrity. Lecture~11 derives these values.

\subsection{From Release 99 to HSPA}
Release~99 carried packet data on dedicated channels with fixed spreading factors and 10--80\,ms
transmission time intervals (TTIs). Allocating a dedicated code to a bursty user wasted code
space, and scheduling in the RNC was slow. HSPA moved intelligence to the NodeB:
\begin{itemize}
  \item \textbf{Release 5 -- HSDPA} \cite{ts25308}: a shared downlink channel (HS-DSCH) using up
        to 15 codes of SF\,16, a 2\,ms TTI, adaptive modulation (QPSK/16QAM), fast hybrid ARQ and
        channel-dependent scheduling in the NodeB, driven by CQI reports from the UE every 2\,ms.
  \item \textbf{Release 6 -- HSUPA} (enhanced uplink): an enhanced dedicated channel (E-DCH)
        with NodeB scheduling and HARQ, up to 5.76\,Mbit/s.
  \item \textbf{Release 7 -- HSPA+}: 64QAM in the downlink, $2\times2$ MIMO, 16QAM uplink, and
        a flat-architecture option.
  \item \textbf{Release 8}: dual-carrier HSDPA (two adjacent 5\,MHz carriers).
\end{itemize}
The philosophical shift matters more than the numbers: the network stopped allocating fixed
resources to users and started scheduling a shared pipe every 2\,ms according to channel
quality -- exactly the paradigm LTE adopts from the start (Lecture~7).

\subsection{IMS: IP multimedia over the PS domain}
The IP Multimedia Subsystem (IMS), introduced in Release~5 \cite{ts23228}, defines how to
offer voice, video and messaging as IP services controlled by SIP \cite{rfc3261}. Its core
elements are the call session control functions: the \textbf{P-CSCF} (proxy, the UE's first
contact point), the \textbf{I-CSCF} (interrogating, entry point of the home network) and the
\textbf{S-CSCF} (serving, which registers the user and executes service logic), plus the HSS,
the evolution of the HLR. In 3G, IMS was rarely used for mass-market voice; its real role came
with LTE, where it is the only way to make calls (VoLTE, Lecture~5).

% ---------------------------------------------------------------------
\section{Mathematical models}

\subsection{Processing gain and the required SINR}\label{sec:pg}
\paragraph{Problem} What signal-to-interference ratio must a WCDMA receiver see before
despreading for a given service quality?
\paragraph{Model} Required energy per bit to noise-plus-interference density
$(E_\mathrm{b}/N_0)_{\mathrm{req}}$, chip rate $W$, service bit rate $R$.
\paragraph{Derivation} Despreading accumulates $W/R$ chips per bit, so
\begin{equation}\label{eq:ebno}
  \frac{E_\mathrm{b}}{N_0} = \frac{W}{R}\cdot\SINR_{\mathrm{wideband}}
  \quad\Longrightarrow\quad
  \SINR_{\mathrm{req}} = \frac{(E_\mathrm{b}/N_0)_{\mathrm{req}}}{G_\mathrm{p}}, \qquad G_\mathrm{p} = \frac{W}{R}.
\end{equation}
\paragraph{Example and interpretation} AMR voice at $R = \qty{12.2}{\kilo\bit\per\second}$:
$G_\mathrm{p} = 3.84\times10^6/12.2\times10^3 = 315 \;(\qty{25.0}{\dB})$. With
$(E_\mathrm{b}/N_0)_{\mathrm{req}} = \qty{5}{\dB}$, the wideband SINR may be as low as
$5-25 = \qty{-20}{\dB}$: the signal is buried 20\,dB below the interference and still decoded.
A \qty{384}{\kilo\bit\per\second} data user has only $G_\mathrm{p} = 10$\,dB and needs
roughly 15\,dB more received power -- which is why high-rate users consume a large share of the
cell's capacity.

\subsection{Uplink load factor, noise rise and pole capacity}\label{sec:load}
\paragraph{Problem} How many simultaneous users can a WCDMA cell carry in the uplink?
\paragraph{Assumptions} $N$ users of the same service, perfect power control, received powers
adjusted so that each user just meets its $E_\mathrm{b}/N_0$; voice activity factor $\nu$;
other-cell interference equal to a fraction $i$ of own-cell interference; thermal noise
power $P_\mathrm{N}$ \cite{holma2010wcdma}.
\paragraph{Derivation} Let $I_{\mathrm{tot}}$ be the total received wideband power and $P_j$ the
received power of user $j$. Its requirement is
\begin{equation}\label{eq:userreq}
  \left(\frac{E_\mathrm{b}}{N_0}\right)_j = \frac{W}{\nu_j R_j}\,\frac{P_j}{I_{\mathrm{tot}}-P_j}
  \quad\Longrightarrow\quad
  P_j = L_j\,I_{\mathrm{tot}}, \qquad
  L_j = \frac{1}{1+\dfrac{W}{(E_\mathrm{b}/N_0)_j\,R_j\,\nu_j}} .
\end{equation}
Summing own-cell users and adding other-cell interference,
$I_{\mathrm{tot}} = (1+i)\sum_j P_j + P_\mathrm{N}$, which defines the uplink load factor
\begin{equation}\label{eq:load}
  \eta_{\mathrm{UL}} = (1+i)\sum_{j=1}^{N} L_j
  \qquad\text{and}\qquad
  \frac{I_{\mathrm{tot}}}{P_\mathrm{N}} = \frac{1}{1-\eta_{\mathrm{UL}}} .
\end{equation}
The ratio $I_{\mathrm{tot}}/P_\mathrm{N}$ is the \emph{noise rise}. It diverges as
$\eta_{\mathrm{UL}}\to1$; the number of users at which this happens is the \emph{pole capacity}
\begin{equation}\label{eq:pole}
  N_{\mathrm{pole}} = \frac{1}{(1+i)L} = \frac{1+\dfrac{W}{(E_\mathrm{b}/N_0)R\nu}}{1+i}.
\end{equation}
\paragraph{Interpretation} \Cref{fig:noise} shows the noise rise: 3\,dB at 50\% load, 6\,dB at 75\%,
10\,dB at 90\%. Each decibel of noise rise is a decibel of lost link budget -- coverage shrinks
as the cell fills up. Operators therefore plan for 50--75\% load, not for the pole.
\paragraph{Limitations} The model assumes perfect power control, a single service and a fixed
$i$; real systems need margins for power-control errors and fast fading.

\begin{acnfigure}
  \acnfig{lecture-03/noise-rise}
  \caption{Uplink noise rise $-10\log_{10}(1-\eta_{\mathrm{UL}})$ from \cref{eq:load}.}
  \label{fig:noise}
\end{acnfigure}

\subsection{HSDPA peak rate}\label{sec:hsdpa}
With $n_{\mathrm{c}}$ codes of spreading factor SF and $b$ bits per symbol, the symbol-level
rate of the HS-DSCH is
\begin{equation}\label{eq:hsdpa}
  R_{\mathrm{peak}} = n_{\mathrm{c}}\,\frac{W}{\mathrm{SF}}\,b\,r_{\mathrm{c}},
\end{equation}
where $r_\mathrm{c}$ is the effective code rate. For $n_\mathrm{c} = 15$, $\mathrm{SF}=16$, 16QAM
($b=4$) and $r_\mathrm{c}\to1$: $15\times240\,\text{ksymbol/s}\times4 = \qty{14.4}{\mega\bit\per\second}$;
the highest standardised UE category reaches \qty{14.0}{\mega\bit\per\second}. 64QAM gives
\qty{21.6}{\mega\bit\per\second} (21.1 in practice), and two MIMO layers or two carriers double
it (\cref{fig:hspa}). These are peak values for one user alone in a cell with excellent SINR;
average cell throughput is several times lower.

\begin{acnfigure}
  \acnfig{lecture-03/hspa-rates}
  \caption{Peak downlink rates along the WCDMA/HSPA evolution (R99 shown with its practical
  384\,kbit/s; categories with the highest standardised rate).}
  \label{fig:hspa}
\end{acnfigure}

% ---------------------------------------------------------------------
\section{Network architecture and diagrams}
\Cref{fig:umts} shows the UMTS network. Compared with GSM (Lecture~1, Figure~3) three things
changed: the RAN uses new nodes and interfaces (NodeB, RNC, Iub, Iur); the core attaches to the
RAN only through Iu, which has a CS and a PS flavour; and the PS user plane is a GTP-U tunnel
that starts at the RNC. The core network itself is recognisably the GSM/GPRS core.

\begin{acnfigure}
  \resizebox{\linewidth}{!}{\acnfig{lecture-03/umts-architecture}}
  \caption{UMTS architecture with the UTRAN, the circuit-switched and the packet-switched
  core domains. Dashed: MAP signalling to the HLR.}
  \label{fig:umts}
\end{acnfigure}

\begin{acnfigure}
  \acnfig{lecture-03/ovsf-tree}
  \caption{OVSF code tree. Assigning $(1,1,1,1)$ blocks its parent $(1,1)$ and the root, but
  its sibling $(1,1,-1,-1)$ remains usable; all codes of one SF are mutually orthogonal.}
  \label{fig:ovsf}
\end{acnfigure}

% ---------------------------------------------------------------------
\section{Step-by-step procedures}
\begin{procedurebox}{PS service request and RAB establishment}\label{proc:rab}
\begin{steps}
  \item The idle UE sends \msg{RRC Connection Request} on the random access channel.
  \item The RNC performs admission control, sets up a radio link in the NodeB (NBAP) and
        answers \msg{RRC Connection Setup}; the UE confirms with \msg{RRC Connection Setup Complete}.
  \item The UE sends a NAS \msg{Service Request} to the SGSN, carried in an RRC
        \msg{Initial Direct Transfer} and the RANAP \msg{Initial UE Message}.
  \item The SGSN may run UMTS AKA; it then sends RANAP \msg{Security Mode Command} with CK and
        IK to the RNC, which starts ciphering and integrity protection towards the UE.
  \item The SGSN requests the bearer with RANAP \msg{RAB Assignment Request}, including QoS and its
        GTP-U tunnel endpoint.
  \item The RNC configures the radio bearer with \msg{RRC Radio Bearer Setup}; the UE answers
        \msg{Radio Bearer Setup Complete}; the RNC returns \msg{RAB Assignment Response} with its own
        tunnel endpoint. User data now flows over the radio bearer and the Iu-PS tunnel.
\end{steps}
\end{procedurebox}

\begin{acnfigure}
  \resizebox{\linewidth}{!}{\acnfig{lecture-03/rab-setup}}
  \caption{RRC connection and RAB set-up for a packet service (\cref{proc:rab}).}
  \label{fig:rab}
\end{acnfigure}

Note the layering: NAS messages (Service Request) travel transparently from UE to SGSN;
access-stratum signalling (RRC) ends in the RNC; RANAP connects the two worlds. This split
between \emph{non-access stratum} and \emph{access stratum} remains in LTE and 5G.

% ---------------------------------------------------------------------
\section{Worked examples}

\begin{examplebox}{Voice pole capacity}\label{ex:pole}
\textbf{Problem.} AMR voice, $R=\qty{12.2}{\kilo\bit\per\second}$,
$(E_\mathrm{b}/N_0)_{\mathrm{req}} = \qty{5}{\dB}$, activity $\nu=0.67$, other-cell interference
ratio $i=0.65$. Find the pole capacity and the number of users at 50\% load.

\textbf{Solution.} $W/((E_\mathrm{b}/N_0)R\nu) = 3.84\times10^6/(3.162\times12\,200\times0.67) = 148.6$.
\Cref{eq:pole}: $N_{\mathrm{pole}} = (1+148.6)/1.65 = 90.6$. Each user loads the cell by
$(1+i)L = 1.65/149.6 = 1.10\%$, so 50\% load corresponds to about 45 users, with a noise
rise of \qty{3}{\dB}.

\textbf{Interpretation.} 45 voice users per 5\,MHz carrier and sector, i.e.\ 9 per MHz -- a
substantial gain over GSM's roughly 94 channels in 10\,MHz with reuse (Lecture~1).
\end{examplebox}

\begin{examplebox}{Mixing services}\label{ex:mix}
\textbf{Problem.} In the cell of \cref{ex:pole}, one \qty{384}{\kilo\bit\per\second} data user
needs $(E_\mathrm{b}/N_0) = \qty{1.5}{\dB}$ with $\nu = 1$. How many voice users does it replace?

\textbf{Solution.} $W/((E_\mathrm{b}/N_0)R\nu) = 3.84\times10^6/(1.413\times384\,000) = 7.08$, so
$L = 1/8.08 = 0.124$ and its load contribution is $1.65\times0.124 = 20.4\%$. One voice user
contributes $1.10\%$: the data user replaces about 18 voice users.

\textbf{Interpretation.} Resource consumption in CDMA scales with $R\cdot E_\mathrm{b}/N_0$, which
is why high-rate packet data on dedicated channels was so costly in Release~99.
\end{examplebox}

\begin{examplebox}{OVSF code budget}\label{ex:ovsf}
\textbf{Problem.} HSDPA reserves 15 codes of SF\,16. What fraction of the downlink code tree is
left for everything else?

\textbf{Solution.} SF\,16 has 16 codes; 15 are used, leaving one SF\,16 branch, i.e.\ $1/16 =
6.25\%$ of the tree, which must hold common channels (pilot, broadcast, paging), the HS-SCCH
control channels and any Release~99 dedicated channels.

\textbf{Interpretation.} In practice operators configure 5--10 HS-PDSCH codes and keep the rest
for voice -- code space, not power, often limits the downlink.
\end{examplebox}

% ---------------------------------------------------------------------
\section{Technology comparison}
\begin{table}[htbp]
\centering\small
\caption{Mobile generations up to 3.5G (indicative peak values).}\label{tab:gen}
\rowcolors{2}{acnmist}{white}
\begin{tabularx}{\linewidth}{@{}l Y Y Y Y@{}}
\toprule
\rowcolor{white} & GSM/GPRS & EDGE & UMTS R99 & HSPA (Rel-5/6) \\
\midrule
Multiple access   & TDMA/FDMA       & TDMA/FDMA         & WCDMA            & WCDMA + shared channel \\
Carrier           & 200\,kHz        & 200\,kHz          & 5\,MHz           & 5\,MHz \\
Peak DL           & 171\,kbit/s     & 474\,kbit/s       & 384\,kbit/s (2\,Mbit/s theor.) & 14\,Mbit/s \\
Handover          & Hard            & Hard              & Soft + hard      & Soft (DCH); serving-cell change (HS-DSCH) \\
Scheduling        & BSC/PCU         & BSC/PCU           & RNC, 10--80\,ms  & NodeB, 2\,ms \\
Core              & CS + PS         & CS + PS           & CS + PS (Iu)     & CS + PS (Iu) \\
Authentication    & One-way (triplet) & One-way         & Mutual AKA (quintet) & Mutual AKA \\
\bottomrule
\end{tabularx}
\end{table}

\paragraph{The CS/PS split and why LTE removed it} Running two core domains means two sets of
nodes, two mobility management procedures (location area and routing area updates), two
charging systems and two sets of interconnection agreements. Every innovation (for example a
new QoS scheme) had to be done twice. Once HSPA proved that a shared packet channel can carry
voice-quality traffic, the economic case for a single, packet-only core became overwhelming:
LTE's EPC (Lecture~9) has no CS domain, and voice is carried by IMS over IP. The price is
complexity elsewhere -- QoS bearers must now guarantee what a circuit guaranteed for free, and
legacy voice requires fallback procedures (CS fallback, SRVCC) during the transition.

% ---------------------------------------------------------------------
\section{Key takeaways}
\begin{takeawaybox}
\begin{itemize}
  \item UMTS = new radio (WCDMA, 5\,MHz, 3.84\,Mchip/s) + evolved GSM/GPRS core, joined by Iu.
  \item The RNC terminates RLC/MAC/RRC, ciphering and integrity; SRNC/DRNC roles enable soft
        handover across RNCs over Iur.
  \item CDMA capacity is interference-limited: $\eta_{\mathrm{UL}} = (1+i)\sum L_j$, noise rise
        $1/(1-\eta_{\mathrm{UL}})$, plan for 50--75\% load.
  \item Processing gain $W/R$ lets low-rate users operate at negative SINR; high-rate users cost
        proportionally more.
  \item HSPA moved scheduling and HARQ to the NodeB with a 2\,ms TTI -- the template for LTE.
  \item Mutual authentication (AKA) and integrity protection fixed the main GSM security gaps.
  \item Two core domains doubled cost and complexity; LTE kept only the packet domain and moved
        voice to IMS.
\end{itemize}
\end{takeawaybox}

% ---------------------------------------------------------------------
\section{Exercises}

\begin{exercisebox}{1}{RNC versus BSC}\label{exr:rnc}
List three functions performed by the UMTS RNC that the GSM BSC did not perform, and explain
which design choice of WCDMA made each necessary.
\end{exercisebox}

\begin{exercisebox}{1}{Why Iur?}\label{exr:iur}
GSM has no interface between BSCs. Explain why UMTS needs Iur, and what would happen to soft
handover at an RNC border without it.
\end{exercisebox}

\begin{exercisebox}{2}{Load budget}\label{exr:load}
Using the parameters of \cref{ex:pole}, how many voice users can be admitted if the noise rise
must not exceed \qty{4}{\dB}? How does the answer change if $i$ drops to 0.4 (better isolated cells)?
\end{exercisebox}

\begin{exercisebox}{2}{HSDPA rate}\label{exr:hsdpa}
A cell configures 10 HS-PDSCH codes. Compute the symbol-level peak rate for QPSK, 16QAM and
64QAM, and the rate with code rate $r_\mathrm{c}=0.75$ for 16QAM.
\end{exercisebox}

\begin{exercisebox}{3}{Cell breathing}\label{exr:breath}
Assume a path-loss exponent of 3.5. By what factor does the cell radius shrink when the load
grows from 50\% to 75\%, if the UE transmit power is the limiting factor? Derive the general
expression for the radius ratio as a function of the two load factors.
\end{exercisebox}

\begin{exercisebox}{3}{Mixed traffic}\label{exr:mixed}
A cell must carry $N_\mathrm{v}$ voice users (parameters of \cref{ex:pole}) and $N_\mathrm{d}$
users at 64\,kbit/s with $E_\mathrm{b}/N_0 = 2$\,dB and $\nu=1$, at a total load of 60\%.
Derive the linear constraint between $N_\mathrm{v}$ and $N_\mathrm{d}$ and plot (or tabulate) the
admissible region.
\end{exercisebox}

\subsection*{Solutions}
\addcontentsline{toc}{subsection}{Solutions}

\begin{solutionbox}{\ref{exr:rnc}}
(1) Macro-diversity combining and soft-handover control -- needed because all cells use the same
frequency. (2) Outer-loop power control (setting the SIR target from frame error rates) --
needed because CDMA is interference-limited and power control is the main resource control.
(3) Ciphering and integrity protection, terminating RLC/MAC/RRC -- moved from BTS/SGSN to keep
security and retransmission in one place despite soft handover across many NodeBs. Admission
control based on load (\cref{eq:load}) is a fourth, because capacity is soft.
\end{solutionbox}

\begin{solutionbox}{\ref{exr:iur}}
In soft handover a UE may be connected to cells of two RNCs. Without Iur the SRNC could not
receive the uplink frames decoded by the other RNC's NodeB, nor send downlink data through it;
the UE would have to perform a hard handover at every RNC border, losing macro-diversity exactly
where links are weakest. Iur lets the DRNC forward frames to the SRNC, which keeps the Iu
connection unchanged until a relocation is worthwhile.
\end{solutionbox}

\begin{solutionbox}{\ref{exr:load}}
Noise rise 4\,dB: $1-\eta = 10^{-0.4} = 0.398$, so $\eta_{\max} = 0.602$. With
$(1+i)L = 1.10\%$, $N = 0.602/0.0110 = 54$ users. With $i = 0.4$: $(1+i)L = 1.4/149.6 = 0.936\%$,
$N = 64$ users -- better cell isolation (sectorisation, antenna tilt) translates directly into
capacity.
\end{solutionbox}

\begin{solutionbox}{\ref{exr:hsdpa}}
Per code: $W/16 = 240$\,ksymbol/s; 10 codes: 2.4\,Msymbol/s. QPSK: 4.8\,Mbit/s; 16QAM:
9.6\,Mbit/s; 64QAM: 14.4\,Mbit/s. 16QAM with $r_\mathrm{c} = 0.75$: $9.6\times0.75 =
7.2$\,Mbit/s.
\end{solutionbox}

\begin{solutionbox}{\ref{exr:breath}}
The extra noise rise between loads $\eta_1$ and $\eta_2$ is
$\Delta = 10\log_{10}\frac{1-\eta_1}{1-\eta_2}$\,dB. With a fixed UE power the tolerable path
loss drops by $\Delta$; for $L_\mathrm{p}\propto d^{n}$, $d_2/d_1 = 10^{-\Delta/(10n)} =
\left(\frac{1-\eta_2}{1-\eta_1}\right)^{1/n}$. For $\eta_1=0.5$, $\eta_2=0.75$: $\Delta = 3.01$\,dB
and $d_2/d_1 = 0.5^{1/3.5} = 0.820$: the radius shrinks by 18\% and the area by 33\%.
\end{solutionbox}

\begin{solutionbox}{\ref{exr:mixed}}
Data user: $W/((E_\mathrm{b}/N_0)R) = 3.84\times10^6/(1.585\times64\,000) = 37.9$,
$L_\mathrm{d} = 1/38.9 = 0.0257$. Constraint:
$1.65\,(N_\mathrm{v}\times0.00668 + N_\mathrm{d}\times0.0257) \le 0.60$, i.e.\
$0.0110\,N_\mathrm{v} + 0.0424\,N_\mathrm{d} \le 0.60$. Corner points: $(N_\mathrm{v},N_\mathrm{d}) =
(54,0)$, $(0,14)$; e.g.\ 30 voice users leave room for 6 data users. One 64\,kbit/s user costs
about 3.9 voice users.
\end{solutionbox}

% ---------------------------------------------------------------------
\section{Further reading}
Holma and Toskala \cite{holma2010wcdma} is the standard reference for WCDMA and HSPA radio
planning, including the load equations of \cref{sec:load}. The UTRAN architecture is in
\cite{ts25401}, RANAP in \cite{ts25413}, HSDPA in \cite{ts25308} and 3G security in
\cite{ts33102}. IMS is specified in \cite{ts23228} on top of SIP \cite{rfc3261}; the general
3GPP architecture is in \cite{ts23002}.

\printbibliography[heading=subbibliography,title={\ACNstrRefs}]
\end{document}
